% ============================================================
%  USPSC-Cap4-Avaliacao_Experimental_4.2-Resultados_por_Abordagem.tex
%  Seção 4.2 — Resultados por Abordagem
% ============================================================

\section{Resultados por Abordagem}
\label{sec:resultados_abordagem}

Esta seção apresenta os resultados individuais de cada abordagem de classificação sobre o conjunto de teste ($n = 373$). Para cada abordagem são reportados: o relatório completo de classificação por classe (\autoref{subsec:metricas_avaliacao}) e a matriz de confusão, que evidencia os padrões sistemáticos de acerto e erro.


% ── 4.2.1 OpLexicon ──────────────────────────────────────────────────────────

\subsection{OpLexicon v3.0}
\label{subsec:res_oplexicon}

O OpLexicon~v3.0 \cite{souza2011oplexicon} classifica cada publicação pela polaridade líquida dos termos presentes no léxico, após a etapa de pré-processamento. O resultado apresentado na \autoref{tab:metrics_oplexicon} revela o desempenho mais baixo entre as quatro abordagens avaliadas.


\begin{table}[!htbp]
	\caption[Métricas de desempenho — OpLexicon v3.0]{Métricas de desempenho do OpLexicon v3.0 ($n = 373$)}
	\label{tab:metrics_oplexicon}
	\centering
	\begin{tabular}{lcccc}
		\hline
		\textbf{Classe}    & \textbf{Precisão} & \textbf{Revocação} & \textbf{F1-Score} & \textbf{Suporte} \\
		\hline		
		Positivo           & 0,18              & 0,68               & 0,29              & 59  \\
		\hline
		Neutro             & 0,81              & 0,13               & 0,22              & 266 \\
		\hline
		Negativo           & 0,20              & 0,46               & 0,28              & 48  \\
		\hline
		\textit{Macro avg} & 0,40              & 0,42               & 0,26              & 373 \\
		\hline
		\textit{Acurácia}  & ---               & ---                & 0,26              & 373 \\
		\hline
	\end{tabular}
	\legend{Fonte: Elaborada pelo autor.}
\end{table}


\begin{table}[!htbp]
	\caption[Matriz de confusão — OpLexicon v3.0]{Matriz de confusão do OpLexicon v3.0. Linhas: rótulo verdadeiro; colunas: rótulo predito.}
	\label{tab:cm_oplexicon}
	\centering
	\begin{tabular}{lccc}
		\hline
		& \textbf{Pred. Positivo} & \textbf{Pred. Neutro} & \textbf{Pred. Negativo} \\
		\hline
		\textbf{Real Positivo} & 40                      & 4                     & 15                      \\
		\hline
		\textbf{Real Neutro}   & 157                     & 34                    & 75                      \\
		\hline
		\textbf{Real Negativo} & 22                      & 4                     & 22                      \\
		\hline
	\end{tabular}
	\legend{Fonte: Elaborada pelo autor.}
\end{table}

O OpLexicon alcançou acurácia de 25,74\% e F1-score macro de 26,12\%, indicando desempenho sistematicamente inferior ao classificador aleatório estratificado. A matriz de confusão (\autoref{tab:cm_oplexicon}) revela o principal padrão de falha: 157 das 266 publicações neutras (59,0\%) foram classificadas como positivo, evidenciando uma tendência de atribuir polaridade positiva a publicações jornalísticas factuais cujo vocabulário financeiro não carrega sinal negativo explícito. A classe neutro, dominante no corpus, obteve revocação de apenas 13\%, o que significa que a abordagem léxica é incapaz de identificar o padrão de ausência de sentimento predominante neste tipo de dado.

% ── 4.2.2 SentiLex-PT ────────────────────────────────────────────────────────

\subsection{SentiLex-PT}
\label{subsec:res_sentilex}

O SentiLex-PT \cite{carvalho2015sentilex} é um léxico de sentimentos para o português construído a partir de entradas morfologicamente anotadas. Os resultados são apresentados na \autoref{tab:metrics_sentilex}.

\begin{table}[!htbp]
	\caption[Métricas de desempenho — SentiLex-PT]{Métricas de desempenho do SentiLex-PT ($n = 373$)}
	\label{tab:metrics_sentilex}
	\centering
	\begin{tabular}{lcccc}
		\hline
		\textbf{Classe}    & \textbf{Precisão} & \textbf{Revocação} & \textbf{F1-Score} & \textbf{Suporte} \\
		\hline
		Positivo           & 0,19              & 0,59               & 0,29              & 59  \\
		\hline
		Neutro             & 0,78              & 0,20               & 0,31              & 266 \\
		\hline
		Negativo           & 0,18              & 0,48               & 0,27              & 48  \\
		\hline
		\textit{Macro avg} & 0,38              & 0,42               & 0,29              & 373 \\
		\hline
		\textit{Acurácia}  & ---               & ---                & 0,29              & 373 \\
		\hline
	\end{tabular}
	\legend{Fonte: Elaborada pelo autor.}
\end{table}

\begin{table}[!htbp]
	\caption[Matriz de confusão — SentiLex-PT]{Matriz de confusão do SentiLex-PT. Linhas: rótulo verdadeiro; colunas: rótulo predito.}
	\label{tab:cm_sentilex}
	\centering
	\begin{tabular}{lccc}
		\hline
		& \textbf{Pred. Positivo} & \textbf{Pred. Neutro} & \textbf{Pred. Negativo} \\
		\hline
		\textbf{Real Positivo} & 35                      & 7                     & 17                      \\
		\hline
		\textbf{Real Neutro}   & 129                     & 52                    & 85                      \\
		\hline
		\textbf{Real Negativo} & 17                      & 8                     & 23                      \\
		\hline
	\end{tabular}
	\legend{Fonte: Elaborada pelo autor.}
\end{table}

O SentiLex-PT obteve acurácia de 29,49\% e F1-score macro de 29,00\%, ligeiramente superior ao OpLexicon, mas ainda muito abaixo de um classificador por maioria de classe (\textit{majority class baseline}). A matriz de confusão (\autoref{tab:cm_sentilex}) evidencia que 129 das 266 publicações neutras (48,5\%) foram classificadas como positivo, padrão análogo ao observado no OpLexicon. A revocação da classe neutra é igualmente baixa (20\%), confirmando a dificuldade estrutural de abordagens léxicas em reconhecer ausência de polaridade em textos de caráter informativo. O SentiLex-PT distribui mais erros entre positivo e negativo do que o OpLexicon, refletindo diferenças na cobertura léxica de cada recurso.

% ── 4.2.3 FinBERT-PT-BR ──────────────────────────────────────────────────────

\subsection{FinBERT-PT-BR}
\label{subsec:res_finbert}

O FinBERT-PT-BR \cite{santos2023finbert} é um modelo BERT especializado para o domínio financeiro em língua portuguesa, pré-treinado sobre corpus de notícias financeiras. Neste trabalho, o modelo é utilizado em modo de inferência direta (\textit{zero-shot} em relação ao corpus local, sem \textit{fine-tuning} sobre as publicações coletadas), precedido apenas da etapa de normalização de espaçamento descrita na \autoref{subsec:fluxo_finbert}, conforme detalhado na \autoref{subsubsec:protocolo_inferencia}. Os resultados são apresentados na \autoref{tab:metrics_finbert}.

\begin{table}[!htbp]
	\caption[Métricas de desempenho — FinBERT-PT-BR]{Métricas de desempenho do FinBERT-PT-BR ($n = 373$)}
	\label{tab:metrics_finbert}
	\centering
	\begin{tabular}{lcccc}
		\hline
		\textbf{Classe}    & \textbf{Precisão} & \textbf{Revocação} & \textbf{F1-Score} & \textbf{Suporte} \\
		\hline
		Positivo           & 0,41              & 0,47               & 0,44              & 59  \\
		\hline
		Neutro             & 0,88              & 0,30               & 0,45              & 266 \\
		\hline
		Negativo           & 0,20              & 0,88               & 0,32              & 48  \\
		\hline
		\textit{Macro avg} & 0,49              & 0,55               & 0,40              & 373 \\
		\hline
		\textit{Acurácia}  & ---               & ---                & 0,40              & 373 \\
		\hline
	\end{tabular}
	\legend{Fonte: Elaborada pelo autor.}
\end{table}

\begin{table}[!htbp]
	\caption[Matriz de confusão — FinBERT-PT-BR]{Matriz de confusão do FinBERT-PT-BR. Linhas: rótulo verdadeiro; colunas: rótulo predito.}
	\label{tab:cm_finbert}
	\centering
	\begin{tabular}{lccc}
		\hline
		& \textbf{Pred. Positivo} & \textbf{Pred. Neutro} & \textbf{Pred. Negativo} \\
		\hline
		\textbf{Real Positivo} & 28                      & 7                     & 24                      \\
		\hline
		\textbf{Real Neutro}   & 39                      & 81                    & 146                     \\
		\hline
		\textbf{Real Negativo} & 2                       & 4                     & 42                      \\
		\hline
	\end{tabular}
	\legend{Fonte: Elaborada pelo autor.}
\end{table}

O FinBERT-PT-BR alcançou acurácia de 40,48\% e F1-score macro de 40,44\%, um resultado superior às abordagens léxicas, mas distante do desempenho documentado pelos autores do modelo sobre o corpus de notícias financeiras de treinamento \cite{santos2023finbert}. O padrão mais relevante na \autoref{tab:cm_finbert} é o viés pronunciado em favor da classe negativo: o modelo classificou 146 das 266 publicações neutras (54,9\%) e 24 das 59 positivas (40,7\%) como negativo, atingindo revocação de 88\% nessa classe à custa de precisão de apenas 20\%. Esse comportamento indica que o modelo generaliza excessivamente o padrão negativo aprendido a partir do corpus de notícias formais --- no qual termos como ``queda'', ``perda'' e ``risco'' são indicadores frequentes de sentimento negativo --- para publicações de formato editorial curto do X/Twitter, nas quais o mesmo vocabulário aparece recorrentemente em contextos factuais e neutros.

% ── 4.2.4 BERTimbau ──────────────────────────────────────────────────────────

\subsection{BERTimbau Ajustado}
\label{subsec:res_bertimbau}

O BERTimbau (\texttt{neuralmind/bert-base-portuguese-cased}) \cite{souza2020} foi submetido a \textit{fine-tuning} sobre a partição de treinamento do corpus rotulado, conforme o protocolo descrito na \autoref{subsubsec:protocolo_inferencia_bertimbau}. Os resultados, apresentados na \autoref{tab:metrics_bertimbau}, revelam desempenho amplamente superior às demais abordagens avaliadas.

\begin{table}[!htbp]
	\caption[Métricas de desempenho — BERTimbau ajustado]{Métricas de desempenho do BERTimbau ajustado ($n = 373$)}
	\label{tab:metrics_bertimbau}
	\centering
	\begin{tabular}{lcccc}
		\hline
		\textbf{Classe}    & \textbf{Precisão} & \textbf{Revocação} & \textbf{F1-Score} & \textbf{Suporte} \\
		\hline
		Positivo           & 0,84              & 0,80               & 0,82              & 59  \\
		\hline
		Neutro             & 0,91              & 0,97               & 0,94              & 266 \\
		\hline
		Negativo           & 0,88              & 0,62               & 0,73              & 48  \\
		\hline
		\textit{Macro avg} & 0,88              & 0,80               & 0,83              & 373 \\
		\hline
		\textit{Acurácia}  & ---               & ---                & 0,90              & 373 \\
		\hline
	\end{tabular}
	\legend{Fonte: Elaborada pelo autor.}
\end{table}


\begin{table}[!htbp]
	\caption[Matriz de confusão — BERTimbau ajustado]{Matriz de confusão do BERTimbau ajustado. Linhas: rótulo verdadeiro; colunas: rótulo predito.}
	\label{tab:cm_bertimbau}
	\centering
	\begin{tabular}{lccc}
		\hline
		& \textbf{Pred. Positivo} & \textbf{Pred. Neutro} & \textbf{Pred. Negativo} \\
		\hline
		\textbf{Real Positivo} & 47                      & 10                    & 2                       \\
		\hline
		\textbf{Real Neutro}   & 7                       & 257                   & 2                       \\
		\hline
		\textbf{Real Negativo} & 2                       & 16                    & 30                      \\
		\hline
	\end{tabular}
	\legend{Fonte: Elaborada pelo autor.}
\end{table}

O BERTimbau ajustado alcançou acurácia de 89,54\% e F1-score macro de 82,84\%, representando ganhos absolutos de 49,1 pontos percentuais em acurácia e 42,4 pontos em F1-score macro em relação ao FinBERT-PT-BR. A matriz de confusão (\autoref{tab:cm_bertimbau}) evidencia que o modelo distribuiu os erros de forma muito mais equilibrada entre as três classes: 257 das 266 instâncias neutras (96,6\%) foram corretamente classificadas, 47 das 59 positivas (79,7\%) e 30 das 48 negativas (62,5\%). Diferentemente do FinBERT-PT-BR, o BERTimbau ajustado não apresenta viés sistemático em favor de nenhuma classe; os erros residuais concentram-se principalmente em confusão entre neutro e as classes minoritárias (10 positivos classificados como neutro e 16 negativos como neutro), um padrão esperado dado o desequilíbrio de classes no corpus.